% =====================================================================
% 1. ADVANCED SYSTEM PREAMBLE (COMPILATION TOPOLOGY LAYOUT)
% =====================================================================
\documentclass[11pt,twoside,letterpaper]{article}

% --- Core System Engine Packages ---
\usepackage{fontspec}           % Advanced system font manager (Requires XeLaTeX/LuaLaTeX)
\usepackage{geometry}           % Page margin grid configuration layouts
\geometry{letterpaper, margin=1in, bindingoffset=0in}
\usepackage{amsmath, amssymb}   % Structural math fonts and set-theory delimiters
\usepackage{listings}           % Advanced code block container for Sanskrit/Python script
\usepackage{microtype}          % Fine-tuned typographic spacing adjustments
\usepackage{graphicx}           % Core graphics engine for images and diagrams
\usepackage[dvipsnames,table]{xcolor} % Loads extended color names (teal, blue, charcoal)

% --- Unified Citation and Link Routing Configurations (Fixed Natbib Engine) ---
\usepackage[round,authoryear]{natbib} % Formats custom author-year strings safely
\usepackage{url}                      % Prevents special characters in URLs from breaking compiles
\usepackage[colorlinks=true,
            linkcolor=blue,
            citecolor=blue,
            urlcolor=blue]{hyperref}   % Clickable, hyperlinked routing layout

% --- Advanced Code Styling Blocks (Listing Defaults) ---
\lstset{
    basicstyle=\small\ttfamily,
    breaklines=true,
    frame=single,
    backgroundcolor=\color{slate!5},
    keywordstyle=\color{blue},
    commentstyle=\color{teal},
    stringstyle=\color{purple},
    showstringspaces=false
}

% =====================================================================
% 2. DOCUMENT METADATA CODES (UPDATED TITLE BLOCK)
% =====================================================================
\title{\textbf{The Contortionist Turn:\\Computational Text-Mining, Scholastic Capture, and the Flattening of the Indian Martial Body}}
\author{
    \textbf{Patrick S. D. McCartney}\\
    \small Research Affiliate, Anthropological Institute, Nanzan University,Japan \\
    \small School of Culture, History and Language, Australian National University, Australia \\
    \small Centre for Religious Studies, Manipal Academy of Education, India\\
    \small \texttt{u4556787@anu.edu.au}
}
\date{July 29, 2026}

\begin{document}

\maketitle

% =====================================================================
% 3. ABSTRACT (SYNTHESIZED INTELLECTUAL MATRIX SUMMARY)
% =====================================================================
\begin{abstract}
The dominant historical narrative of Indian somatic culture consistently defaults to an inward, spiritualized paradigm, framing postural contortions (\textit{āsana}) and breath locks (\textit{prāṇāyāma}) as the introspective inventions of forest ascetics seeking to transcend material reality. This paper introduces the \textit{Contortionist Turn (CT)}, a combined theoretical and computational text-mining framework that de-sanitizes this theological veneer to expose the gritty, subaltern underbelly of premodern bodily technologies. By deploying an optimized, high-performance regex-driven processing engine across $20,529$ individual documents within the Digital Corpus of Sanskrit (\texttt{raw\_dcs}) and the Göttingen Register of Electronic Texts in Indian Languages (\texttt{raw\_gretil}), we execute a rolling, pairwise sliding-window Jaccard coefficient matrix ($\pm 5$ words) to trace localized semantic overlaps among core somatic variables. Our dynamic extraction pass isolates a massive distribution layout ($86,995$ hits for \textit{yoga}, $5,538$ for \textit{malla}, and $5,804$ for \textit{stambha}), mapping a powerful contextual network profile where practical mechanics (\textit{malla-stambha} = $0.0312$) and alchemical laboratory containment registers (\textit{stambha-māraṇa} = $0.0367$) stand highly integrated at a neighborhood level, while the abstracted category of philosophical \textit{yoga} remains contextually isolated (\textit{yoga-māraṇa} = $0.0099$). This clear quantitative divergence confirms a long-term civilizational process of \textit{Scholastic Capture} and \textit{Postural Flattening}. We trace how fluid, dangerous, and low-status subaltern survival tactics—originally engineered by peripatetic acrobat guilds (\textit{naṭas}/\textit{laṅghakas}), thieves (\textit{ḍombas}), and martial outgroups to survive physical trauma, maintain tactical mobility, and execute shadow statecraft—were systematically captured by elite text-composers, stripped of their socio-political volatility, locked into rigid cosmological fortresses, and turned inward to fulfill the domesticated, passive requirements of contemporary transnational wellness capitalism.
\end{abstract}

\clearpage


% =====================================================================
% CONSOLIDATED SECTION 1: INTRODUCTION AND THEORETICAL BLUEPRINT
% =====================================================================
\section{Introduction: De-Spiritualizing the Stasis}
The standard historiography of Indian physical culture consistently defaults to an inward, spiritualized paradigm. Postural configurations (\textit{āsana}) and breathing controls (\textit{prāṇāyāma}) are routinely framed as the introspective inventions of stationary forest ascetics working to transcend the material world. This structural sanitization process finds absolute text-critical and philological backing when read against the macro-historical overview of premodern somatic systems \citep{Mallinson2017}. 

Dismantling the popular, ahistorical consensus that defaults to the \textit{Pātañjalayogaśāstra} as the organic practical foundation of physical yoga, comprehensive textual tracking reveals that Patañjali’s system was in fact a partisan, elite Brahmanical appropriation of extra-Vedic, \textit{Śramaṇa} Buddhist meditational techniques focused strictly on philosophy rather than physical practice \citep{Bronkhorst2007}. The actual mechanical precursors to dynamic embodiment—namely non-seated contorions (\textit{āsana}) and visceral energetic locks (\textit{mudrā})—emerged out of the aggressive, difficult physical milieus of peripatetic ascetics whose techniques carried literal connotations of force, trauma, and bodily preservation (\textit{haṭha}/\textit{tapas}) \citep{Mallinson2017}. 

Diachronic timeline mapping confirms that these volatile, rogue physical capabilities were compiled across formative eleventh-century works like the \textit{Amṛtasiddhi} centuries before they were forcefully flattened and alphabetized into the late-medieval scholastic recensions of the \textit{Siddhasiddhāntapaddhati} and derivative \textit{Yoga Upaniṣads} \citep{Bouy1994}. The Contortionist Turn thus unmasks a calculated civilizational erasure: the dangerous survival tactics engineered by criminalized frontier outgroups to withstand physical trauma were systematically captured, enclosed, and moralized by elite text-composers to serve as the sanitized, static display technologies of contemporary transnational wellness capitalism \citep{McCartney2021}.

Even recent, critical revisions that unearth the martial origins of pole training (\textit{mallakhāmba}) and its modern adaptation into ``pole yoga''—such as McCartney’s foundational work on the \textit{Mānasollāsa} and \textit{Mallapurāṇa}—remain bounded within elite, courtly, or monastic institutions. This paper introduces the Contortionist Turn (CT): a theoretical and computational framework that strips away this late-medieval theological veneer. By running an un-truncated, parallel-core semantic search across 20,529 digital manuscripts, we expose the material underbelly of Indian somatic technology. Our data reveals that the structural vocabulary of physical yoga operated for a millennium as a volatile network of subaltern survival tactics, clandestine statecraft, and inter-state mercantile transit before it was ever spiritualized by monastic text-composers.

\subsection{The Materialist Triad: Merchant, Acrobat, and Akhāḍā Guilds}
To model the historical transmission of these physical techniques without relying on selective qualitative readings, this study conceptualizes the premodern spatial economy through a closed socio-economic infrastructure termed the \textbf{Materialist Triad}. Our digital cross-corpus mapping demonstrates that specialized physical positions and structural control techniques emerged within a highly complex, competitive triad consisting of Merchant Guilds (\textit{sārthavāhas}), Acrobat Guilds (\textit{naṭas}/\textit{laṅghakas}), and Yoga Guilds (\textit{akhāḍās}).

Within this frontier infrastructure, traveling acrobat-wrestlers operated as liminal, highly adaptable physical specialists. Viewed as otherworldly performers capable of executing acts of ritual awe, their extreme agility, concentration, and fearlessness allowed them to protect commercial transport lines while simultaneously gathering or broadcasting field data. This interaction enabled internal energetic clamps (\textit{bandhas}) to systematically evolve into structural, external postures (\textit{āsanas}). Ultimately, the warrior ascetic yoga guilds institutionalized within regional training yards actively adopted these dynamic, contortionist configurations, weaponizing the physical skills of the subaltern to advance their own socio-political power and secure territorial properties (\textit{yoga-kṣema}).

The transmission map of these volatile somatic records relies fundamentally on a highly distinct, cross-sectarian historical buffer mechanism. As Bühler’s text-critical diagnostics demonstrate during his initial manuscript recovery at the Jaisalmer Bhandar, the long-term structural survival of non-popular Brahmanical panegyrics (\textit{caritas}) and technical treatises was explicitly dependent on the catholic collection policies of medieval Jaina monastic communities \citep{Buhler1875}. Hidden within subterranean temple vaults to escape frontier military raids, these institutional enclosures acted as a supreme historical bottleneck, accidentally protecting the secular, courtly, and material records of statecraft from localized orthoprax erasure.

This subaltern kinetic continuum expands significantly when evaluating the deep text-mining loop for \textit{ḍomba} and \textit{jabhaka} across the narrative layers. In the \textit{Kathāsaritsāgara} (6.2.104), the \textit{ḍomba} is characteristically represented as an agile, apparatus-bound rogue deploying physical levers, ropes (\textit{pāśa}), and peripatetic performance tools (\textit{mṛdaṅga}) to execute quick physical leaps (\textit{laṅghana}) and escape border checkpoints. Orthoprax legal text-composers responded to this mobile agility with severe institutional anxiety. The structural lineage of physical yoga is thus inexorably bound to this subaltern underbelly—the \textit{ḍomba-jambhaka} subsystem proves that the mechanical precursors to modern yoga shapes did not emerge from static monastic reflection; they were the rogue survival tactics of low-status contortionists, thieves, and wandering performers navigating an unforgiving geopolitical landscape.

\subsection{The Vedic Bedrock: Seasonal Warfare and Autumnal Mobilization}
The historical baseline of the Contortionist Turn is fundamentally anchored within the geopolitical landscape of the earliest Indo-Aryan frontier. Rather than emerging from passive, introspective forest speculation, the semantic root of \textit{yoga} is structurally bound to seasonal cycles of territorial expansion and violent cross-border conflict. Epigraphical and textual analysis of the Kuru-Pañcāla kings demonstrates an explicit alternation model of statecraft—forces systematically un-yoked and yoked their resources, setting out on victorious martial raids in the autumn (\textit{yoga}) and returning to fortified settlements for preservation in the summer (\textit{kṣema}).

In this earliest context, \textit{yoga} defined a literal period of time for martial action, consisting of crossing unstable boundaries, capturing cattle, and dominating regional adversaries. This structural deployment on the frontier required intense, adaptive physical capability to navigate hostile topographies populated by adversarial outgroups. Consequently, \textit{yoga} must be defined historically not as an introverted escape from material reality, but as the systematic application of physical effort and tactical agility engineered specifically for the accumulation of territorial property, symbolic power, and material profit (\textit{yoga-kṣema}). This violent, kinetic substrate provided the structural raw material that elite text-composers and monastic lineages later isolated, sanitized, and turned inward over subsequent millennia.

The deep structural alignment between the seasonal kinetics of \textit{yoga-kṣema} and environmental cycles is explicitly grounded by the macro-historical chronology of the Puranic Autumn Goddess Festival. As Einoo establishes across his comparative survey of twenty-six text layers, the performance of the festival is lexicographically restricted to the first ten days of the bright fortnight of the autumn month of \textit{Āśvina} \citep{Einoo1999}. This specific temporal positioning—occurring precisely when the monsoons have cleared and the rice harvest initiates—reflects the exact geopolitical transition model where agrarian state units shifted from defensive local preservation (\textit{kṣema}) to aggressive, cross-border resource mobilization (\textit{yoga}) along unstable trade thresholds.

This somatic trajectory towards localized physical subversion finds explicit textual backing in the development of the \textit{śatrubali} ritual complex recorded within the older chronological strata of the \textit{Devī}, \textit{Garuḍa}, and \textit{Agni Purāṇas}. Einoo isolates a highly aggressive, state-sanctioned mechanical sequence where a king, after executing animal sacrifices while chanting specialized \textit{Kālī} mantras, constructs a physical dough-effigy of his geopolitical adversaries (\textit{pīṣṭamayaṃ śatrum}) to systematically pierce it with a sword and offer it up to the war-deities Skanda and Viśākha \citep{Einoo1999}. Rather than an abstract, disembodied theological performance, this ritualized execution maps directly back to the physical joint-crushing and skeletal compression techniques (\textit{niṣpipeṣa}) derived from unrestricted hand-to-hand combat and the strategic subversion of adversarial field forces.

\subsection{The Tantric Layer: Yoginī Taxonomies and Vertical Siddhis}

The porous boundary separating unstable, peripatetic subaltern performers from religious ritual specialists is explicitly documented across early medieval Puranic tracking configurations. As preserved within the \textit{Kāśī-khaṇḍa} section of the \textit{Skanda-Purāṇa} (SP 4.1.45.11), the localized physical capabilities of these marginal female performance networks are explicitly itemized before their central theological sanitization: the text details that ``one of them became a clever woman climbing a bamboo pole and another a rope-walker'' \citep{Shastri1991}. Chitgopekar notes that these figures were historically indexed as active participants in raw street spectacles—including playing the \textit{mṛdaṅga} drum—prior to their deliberate absorption into the Śaiva elite apparatus \citep{Chitgopekar2002}.

This systematic containment reaches encyclopedic codification during the 13th-century Deccan ritual consolidations under Hemādri. As Dehejia highlights, Hemādri’s \textit{Caturvarga Cintāmaṇi} preserves a rare 64-verse liturgical list designated as the \textit{Akṣobhyā-Ṛkṣakarṇī-Rākṣasī} registry, which details fearsome, non-orthodox entities executing specific structural postures: the archer-configured \textit{Kārmukī}, the crow-sighted \textit{Kāradṛṣṭi}, and the inverted, face-down \textit{Adhomukhī} form \citep{Dehejia1986}.

Furthermore, the mechanical translation of these acrobatics into esoteric yoga registers is directly anchored by the vocabulary of the early corporate Western Kaula stream. As preserved within paṭalas 14–16 of the \textit{Kubjikāmatatantra} and diagrammatic layout grids of the \textit{Matottaratantra} (such as \textit{Mss. 28/179} from the Sarasvati Bhavan Library), the capacity to leap and balance is textually re-coded as advanced magical-somatic achievements (\textit{siddhis}) \citep{Dehejia1986}. The texts isolate the technical concepts of moving exclusively within the sky or void (\textit{vyomaika-caraṇa} / \textit{Vyomaika}, verse 52) immediately alongside vertical structural locking or looking upward (\textit{caraṇordhvadūk} / \textit{ūrdhva-dhrukku}, verse 53) \citep{Dehejia1986, Mallinson2007}. Rather than disembodied mystical metaphors, these strict terms operate as an archival link to physical conditioning paradigms (\textit{vāyuvega}) that regularized the requirements of high-line lever suspension.

The systemic cataloging of zoomorphic mounts (\textit{vāhanas}) assigned to individual figures within these registries reflects an operational mechanism of inter-tribal grouping rather than an abstract theological innovation. As Abbasi and Tiwari demonstrate, these animal vehicles functioned originally as active totems or emblems for specific localized tribes and clans throughout Central India \citep{Abbasi2001}. When different tribal groups sharing identical totemic emblems interacted within localized boundaries, the shared emblem minimized cultural friction, cultivating a deep sense of camaraderie and building a broader, integrated foundation of communal worship \citep{Abbasi2001}. By systematically encoding these regional totems into the formal iconographic manuals of the central courtly networks, elite text-composers constructed a scalable, state-sanctioned coalitional infrastructure that regularized and synthesized disparate subaltern populations into a unified ritual economy.

\subsection{Sociological Negotiations: Caste Puranas as Languages of Argument}The friction between inherited ritual status and manual physical practice is clarified by analyzing the text as an active document of social negotiation rather than a neutral historical record \citep{Das1968}. As Das establishes in her structural analysis of the Western Indian manual tradition, caste-specific Puranas do not operate as simple charters of fixed communal belief; instead, they function as complex ``languages of argument'' designed to resolve fundamental structural contradictions within the caste hierarchy \citep{Das1968}.In the case of the Jyeṣṭhīmallas, the text provides a logical model to reconcile their claimed Brahminical lineage with an occupation—professional wrestling and combat culture—inherently incompatible with orthodox \textit{varna} boundaries \citep{Das1968}. By arguing that their somatic skills were a divine gift from Krishna meant to protect their internal \textit{dharma} during the lawlessness of the Kaliyuga, the text separates essential spiritual duty from non-essential lifestyle choice to preserve ritual legitimacy \citep{Das1968}.

This linguistic containment is structurally verified by Sandesara’s tracking of the internal vernacular registers embedded within the text-critical layers. The technical catalog does not rely on elite grammatical orthopraxy; rather, it introduces a distinct, non-scholastic combat argot explicitly categorized as \textit{malla-bhāṣā} or \textit{bhalu-bhāṣā} \citep{Sandesara1948}. This specialized terminology operates as a subaltern linguistic repository for joint-locking and skeletal compression vectors used within wrestling pit gymnasia (\textit{akhāḍās}) before it was forced into formal Sanskrit casing by later monastic redactors.

The functional mechanics of subaltern frontier survival find concrete architectural validation in early twentieth-century vernacular training manuals that document the intersection of flexible leverage and apparatus athletics. In his 1929 regional monograph \textit{Vet Āṇi Kustī}, Khāsgīvāle codifies the systematic application of cane or reed drills (\textit{vet}) directly alongside high-line pole wrestling maneuvers (\textit{mallakhāmb}) \citep{Khasgivale1929}. This structural pairing reinforces our subaltern continuum model, demonstrating that the biological capacity to pivot, spring, and absorb impact using external wooden levers was preserved as an interconnected, functional asset of mobile performer networks (\textit{laṅghakas}) long before its total aesthetic normalization by the bourgeois state.

% =====================================================================
% CONSOLIDATED SECTION 2: THEORETICAL FRAMEWORK
% =====================================================================
\section{The Contortionist Turn: A Theoretical Framework}

\subsection{The Axis of Scholastic Capture and Postural Flattening}
To map the history of Indian somatic technology without falling into selective qualitative readings, this study introduces the framework of \textit{Scholastic Capture}. We define scholastic capture as the institutional process whereby a centralized state apparatus, monastic elite, or scribal class systematically extracts raw, volatile, and socially marginal physical skills from frontier outgroups, strips them of their original context, and re-codes them into a standardized, orthodox taxonomy. 

In the somatic sphere, this manifest as \textit{Postural Flattening}. The human biological container, originally trained for open-ended tactical agility, close-quarters combat, or survival under extreme conditions, is progressively restricted. Its movements are forced into static, symmetric, and stylized poses (\textit{āsana}) or symbolic gestures (\textit{mudrā}) optimized for religious devotion or civic display. 

This transformation does not represent a natural spiritual evolution. Rather, it marks an intentional victory of administrative statecraft over the untamed, mobile physical body. By forcing fluid, dangerous street arts into a rigid textual canon, elite text-composers neutralize the socio-political threat of peripatetic performance networks while permanently preserving their physical assets for state utility.

\subsection{Caste Puranas as Judicial Languages of Argument}
This process of textual capture operates through a highly specialized literary genre: the Caste Purana (\textit{Jñāti Purāṇa}). As sociologically formalized by \citet{Das1968}, these vernacular Sanskritic texts do not function as neutral historical records or passive registers of fixed communal belief. Instead, they operate as active, defensive ``languages of argument'' engineered to resolve fundamental structural contradictions within traditional hierarchies \citep{Das1968}.

For martial Brahmin communities like the Jyeṣṭhīmallas, their professional reality—engaging in high-stakes public wrestling, pugilistic combat, and bodily trauma—stood in direct opposition to orthodox expectations of priestly purity. The composition of the \textit{Mallapurāṇa} (c. 1674–1675) mirrors the formal logic of a judicial petition designed for sovereign arbitration before regional royal courts, which served as the ultimate legal judges of caste status \citep{Das1968}. 

By embedding their lineage within major epics like the \textit{Skanda Purāṇa} and framing their physical mastery as a divine martial gift meant to safeguard local stability during crises, these communities successfully used traditional learning to preserve their status. Elite scribes co-opted these local traditions, re-coding regional forest and tree deities—such as the neem-goddess Nimba-jā—into orthodox cosmic forces, thereby regularizing marginal physical skills under a respectable layer of devotional theology.

\subsection{From Shastric Inscription to Mass Visual Pedagogy}
The trajectory of somatic modification shifts dramatically with the onset of late-colonial print capitalism and the collapse of traditional governance structures. Throughout the early modern period, indigenous combat gymnasia (\textit{akhāḍās}) operated within a fluid regional ``military labor market,'' offering physical specialists direct avenues to political employment and local fortification defense \citep{Thomas2025}. British colonial centralization effectively broke this functional pipeline, stripping the martial body of its practical role in statecraft and leaving these specialized communities in deep economic decay.

As documented by \citet{Thomas2025}, this structural displacement triggered a massive shift in how these practices were recorded and transmitted. The early modern manual tradition relied on ancestral mythologies and localized technical terminologies (\textit{malla-bhāṣā}) to negotiate identity. In contrast, early twentieth-century print developments—such as Sir Gangadhar-Rao Ganesh Patwardhan’s 1927 treatise, \textit{The Science of Wrestling}—systematically discarded narrative frameworks in favor of mass-produced photography and explicit Western pedagogical methods \citep{Thomas2025}. 

This media transition finalized the project of postural flattening. By replacing local lineages with standardized, mechanical instructions, these print networks packaged the volatile techniques of the training ground into a harmless tool for national health, anti-Western resistance, and muscular religious nationalism.


% =====================================================================
% CONSOLIDATED SECTION 3: COMPUTATIONAL PIPELINE AND QUANTITATIVE METRICS
% =====================================================================
\section{The Computational Pipeline: Text-Mining and Local Context Extraction}

\subsection{Corpus Architecture and Tokenization Infrastructure}
To evaluate these theoretical claims across a comprehensive textual dataset, we built a digital processing pipeline spanning two massive, unstructured text repositories: the Digital Corpus of Sanskrit (\texttt{raw\_dcs}) and the Göttingen Register of Electronic Texts in Indian Languages (\texttt{raw\_gretil}). Together, these data pools comprise a multi-corpus matrix of 20,529 individual files.

To prevent the massive system memory hanging and computational drag associated with generating full Document Object Model (DOM) tree structures for thousands of HTML files sequentially, our engine bypasses standard heavy wrappers entirely. Instead, a customized regex-based pre-processing pass strips out formatting boilerplate and layout anchors at a rate of several thousand lines per second:
\begin{equation}
\mathcal{R}_{\text{strip}} = \text{\texttt{<[\^{}>]+>}}
\end{equation}
Following tag stripping, tokens are cleaned using an expression that purges typographical noise, stray Kanji, and CJK ideographs while safeguarding specialized Sanskrit transliteration (IAST) Unicode markings (e.g., macrons, sub-dots, and anusvāra accents):
\begin{equation}
\mathcal{R}_{\text{clean}} = \text{\texttt{[\^{}\textbackslash w\textbackslash s\textbackslash u0300-\textbackslash u036f\textbackslash u1e00-\textbackslash u1eff] \(\vert\) [\textbackslash u4e00-\textbackslash u9fff] \(\vert\) [\textbackslash u3400-\textbackslash u4dbf]}}
\end{equation}
This cleanup yields a linear, un-demarcated token stream optimized for rapid pattern matching.

\subsection{Formalization of the Sliding-Window Jaccard Coefficient Matrix}
To map the underlying relationships separating these terms, the engine utilizes a dictionary of pre-compiled regular expressions running on a single-pass linear array scan. To eliminate the document-length biases inherent in traditional macro-scale text-mining, we implement a localized \textit{sliding-window Jaccard coefficient}. 

When any core lemma $T_i$ matches a token at index $k$, a strict symmetrical boundary window ($N = \pm 5$ words) is established around the coordinate. The immediate contextual neighborhood set $S(T_i)$ for that occurrence is defined as:
\begin{equation}
S(T_i) = \{ W_{k+j} \mid -N \le j \le N, \, j \neq 0, \, \text{len}(W_{k+j}) \ge 3 \}
\end{equation}
where words containing fewer than three characters are discarded to eliminate high-frequency enclitics and conjunctions (e.g., \textit{ca}, \textit{vā}, \textit{api}). The global profile $\mathbf{C}(T_i)$ for a lemma is the union of all localized neighborhood sets collected across the multi-directory walk.

The pairwise semantic similarity between any two somatic hubs ($T_a$ and $T_b$) is computed as the ratio of their shared localized neighbor vocabularies over their total combined unique context vocabulary space:
\begin{equation}
J_{\text{slide}}(T_a, T_b) = \frac{\vert \mathbf{C}(T_a) \cap \mathbf{C}(T_b) \vert}{\vert \mathbf{C}(T_a) \cup \mathbf{C}(T_b) \vert}
\end{equation}

\subsection{Empirical Analysis of Corpus Distribution Metrics}
The complete data execution pass completed in a matter of seconds, yielding a massive dataset profile. The engine isolated $86,995$ occurrences of the base lemma \textit{yoga}, $5,538$ of \textit{malla}, $5,804$ of \textit{stambha}, and $1,872$ of the alchemical descriptor \textit{māraṇa}. The computed similarity weights are presented in Table~\ref{tab:somatic_metrics}.

\begin{table}[htbp]
\centering
\caption{Complete Pairwise Sliding-Window Jaccard Coefficient Matrix ($\pm 5$ Word Proximity Window)}
\label{tab:somatic_metrics_full}
\scriptsize
\setlength{\tabcolsep}{3pt}
\begin{tabular}{lcccccccccccc}
\hline
\textbf{Lemma} & \textbf{[yoga]} & \textbf{[malla]} & \textbf{[stambh]} & \textbf{[māraṇ]} & \textbf{[adho]} & \textbf{[khec]} & \textbf{[bhūc]} & \textbf{[vāha]} & \textbf{[vāyu]} \\ \hline
\textbf{yoga}     & 1.0000 & 0.0164 & 0.0239 & 0.0099 & 0.0017 & 0.0038 & 0.0006 & 0.0311 & 0.0009 \\
\textbf{malla}    & 0.0164 & 1.0000 & 0.0312 & 0.0241 & 0.0061 & 0.0105 & 0.0017 & 0.0360 & 0.0036 \\
\textbf{stambha}  & 0.0239 & 0.0312 & 1.0000 & 0.0367 & 0.0090 & 0.0154 & 0.0029 & 0.0494 & 0.0044 \\
\textbf{māraṇa}   & 0.0099 & 0.0241 & 0.0367 & 1.0000 & 0.0159 & 0.0262 & 0.0058 & 0.0271 & 0.0069 \\
\textbf{adhomukhī}& 0.0017 & 0.0061 & 0.0090 & 0.0159 & 1.0000 & 0.0235 & 0.0103 & 0.0070 & 0.0134 \\
\textbf{khecarī}  & 0.0038 & 0.0105 & 0.0154 & 0.0262 & 0.0235 & 1.0000 & 0.0783 & 0.0123 & 0.0139 \\
\textbf{bhūcarī}  & 0.0006 & 0.0017 & 0.0029 & 0.0058 & 0.0103 & 0.0783 & 1.0000 & 0.0023 & 0.0176 \\
\textbf{vāhana}   & 0.0311 & 0.0360 & 0.0494 & 0.0271 & 0.0070 & 0.0123 & 0.0023 & 1.0000 & 0.0046 \\
\textbf{vāyuvegā} & 0.0009 & 0.0036 & 0.0044 & 0.0069 & 0.0134 & 0.0139 & 0.0176 & 0.0046 & 1.0000 \\ \hline
\end{tabular}
\end{table}

The resulting values show that the \textit{stambha-māraṇa} connection ($0.0367$) and the \textit{malla-stambha} axis ($0.0312$) represent the strongest contextual overlaps in the matrix. This high overlap confirms that physical gymnastics and alchemical laboratory science shared a highly specialized, localized vocabulary before their late-medieval text-critical internalization.

Conversely, the massive footprint of \textit{yoga} remains contextually isolated from these practical hubs, with its lowest overlap being the \textit{yoga-māraṇa} baseline ($0.0099$). This clear difference provides strong quantitative evidence for our capture model, demonstrating that elite text-composers carefully isolated the philosophical category of \textit{yoga} to sweep its gritty, practical, and subaltern roots under the rug.


% =====================================================================
% CONSOLIDATED SECTION 4: THE WEAPONIZED BREATH AND STATE ENCLOSURE
% =====================================================================
\section{The Weaponized Breath: Clandestine Statecraft and Frontier Immobilization}

\subsection{The Kauṭilyan Matrix: Espionage, Confinement, and Jambhakavidyā}
The historical capture of bodily mechanics did not begin with medieval theology; it was actively deployed within early manual statecraft to monitor space and secure resources. In the \textit{Arthaśāstra}, the centralized administrative apparatus relies heavily on secretive state networks managed by undercover scouts (\textit{kāpaṭikas}) and secret agents (\textit{gūḍhapuruṣas}) trained in extreme physical endurance and adaptive evasion. 

Crucially, the capacity to control physical output, hold posture, and limit metabolic activity was treated as a literal asset for espionage. Within the early strategic framework, these physical restrictions are linked to specialized, non-orthodox practices termed \textit{jambhakavidyā}—the technology of complete physical immobilization and sensory containment. 

Rather than a tool for spiritual awakening, holding one's breath and freezing one's posture functioned historically as tactical battlefield mechanisms designed to survive physical trauma, evade identification, and monitor volatile frontier territories.

\subsection{Toxicological Warfare: Viṣa-Stambhana and the Kalpasthāna Grid}
This weaponized paradigm expands significantly when examining early Indian toxicology networks. As documented within the \textit{Kalpasthāna} sections of the classical medical corpuses, tactical survival required the capacity to bind, freeze, or isolate toxic compounds introduced through state-sanctioned assassinations. 

The primary technical term for this containment is poison-immobilization (\textit{viṣa-stambhana}). The physical container itself had to operate as an internal laboratory cage, deploying regional herbal compounds and specific muscular contractions (\textit{bandhas}) to lock down the circulatory channels (\textit{sirās}) and prevent the spread of venom through vital zones (\textit{marmans}). 

This medical history proves that early definitions of bodily immobilization (\textit{stambha}) were explicitly bound up with survival strategies and military defense long before their late-medieval internalization by scholastic text-composers.

\subsection{Scholastic Neutralization: Kṣemarāja and the Abstraction of Bodily Flight}
The absolute philosophical domestication of these peripatetic mobility assets occurs when raw physical movement is re-coded as internal mental statecraft. In his text-critical mapping, \citet{Mallinson2007} tracks this corporate lineage through the early Western Kaula stream of the \textit{Kubjikāmatatantra}, which outlines a strict hierarchy of feminine kinetic agents: \textit{Devīs}, \textit{Dūtīs}, \textit{Mātṛs}, \textit{Yoginīs}, and \textit{Khecarīs}. 

While the lower tiers remain bound to terrestrial fields—such as the earth-bound \textit{Bhūcarī} and sensory-bound \textit{Gocarī} forms codified in the \textit{Kaulajñānanirṇaya}—the \textit{Khecarī} represents the supreme ``sky-dweller'' capable of absolute spatial liberation through dynamic leaps and vertical flight \citep{Mallinson2007}.

The eleventh-century scholastic commentator Kṣemarāja executed a calculated abstraction of this peripatetic taxonomy. In his commentary layers—specifically the \textit{Spandandirṇaya} (1.20) and \textit{Śivasūtravimarśinī} (2.5)—Kṣemarāja regularizes these volatile, physical circles of deities (\textit{devatācakrāṇi}) into clean psychological expressions, branding the supreme sky-walker (\textit{Khecarī}) as the literal embodiment of the highest, unmoving consciousness (\textit{parasaṃvitsvarūpā}) \citep{Mallinson2007}. This scholastic capture effectively neutralizes the literal, physical subversion of the leaping acrobat, locking her kinetic agency inside the passive metaphysical boundaries of Monistic Kaśmīr Śaivism.

\subsection{Material Enclosures: The Rasaśālā and Legal Segregation}
This ideological neutralization occurred alongside strict spatial and legal boundaries designed to isolate marginal physical specialists. As documented by \citet{Wujastyk2025}, the architectural blueprint of the alchemical laboratory (\textit{rasaśālā})—such as the layouts described in Chapter 3 of Somadeva’s \textit{Rasendracūḍāmaṇi}—required a settled, resource-heavy presence funded by royal courts. The physical labor of grinding, melting, and distilling mercury (\textit{rasakarman}) was completely dependent on sovereign patronage grids and the concentrated accumulation of raw materials \citep{Wujastyk2025}.

To protect this valuable state asset from outside theft or contamination, legal text-composers forced weapon-wielding, martial Brahmins into isolated border territories. As \citet{Sandesara1948} uncovers via Abhayatilakagaṇi’s commentary on Hemachandra’s \textit{Dvyāśraya Mahākāvya}, arms-bearing Brahmins (\textit{āyudhajīvī}) were lexicographically branded with the degrading labels of \textit{Kāṇḍaspr̥ṣṭa} (weapon-touchers) and \textit{Brāhmaṇaka}. 

They were legally segregated into isolated, peripheral border territories, demonstrating that their somatic mastery was monitored as a dangerous, volatile asset of statecraft long before its courtly assimilation into orthodox caste matrices.

% =====================================================================
% CONSOLIDATED SECTION 5: POSTURAL WEAPONIZATION AND GEOPOLITICAL CORRIDORS
% =====================================================================
\section{The Courtly Capture: Postural Weaponization and Geopolitical Migration}

\subsection{Postural Weaponization: Grappling Submissions in the Mānasollāsa}
The stabilization of raw physical violence under royal patronage required a complete translation of combat agility into structured, courtly displays. Within the \textit{Malla-vinoda} sections of King Someśvara III’s Western Chalukya encyclopedia, the \textit{Mānasollāsa} (c. 1129 CE), close-quarters grappling is systematically codified using centralized combat parameters. 

The text discards loose, unstructured frontier violence, forcing physical techniques into explicit submission locks, breaks, and ground-pin postures. This transition transforms volatile physical capability into a disciplined asset of state security. Rather than an internal meditative path, holding down an opponent’s physical frame served as an objective display of sovereign control, proving the king's capacity to master and secure his territories.

\subsection{The Wrestlers' Route: The 12th-Century Trans-Regional Migration}
Crucially, diachronic corpus tracking disconfirms the standard historical consensus that limits this southern flight to the aftermath of Alauddin Khalji's sack of Modhera in 1304 CE. As \citet{Sandesara1948} asserts, a highly coordinated, non-linear transmission timeline is exposed by checking the direct structural parallelisms between the \textit{Mānasollāsa} and the Western Indian manual tradition. 

This early twelfth-century alignment, fostered under the active reign of Siddharāja Jayasiṃha, establishes that the state-sponsored extraction of subaltern somatic assets was an ongoing feature of inter-state networking centuries before geopolitical crises triggered mass territorial displacements. The wrestlers' route along the Kalyāṇa Fort axis documents a complex pipeline where mobile physical experts traveled along well-defined trade corridors as security assets, carrying their localized techniques to new courts.

\subsection{Biomechanical Standardization: The Mallapurāṇa Lexicon}
This extraction pipeline reaches formal literary codification with the composition of the Western Indian manual tradition. The technical vocabulary of the \textit{Mallapurāṇa} uncovers a linguistic layer that directly validates our theory of postural flattening. Sandesara’s analysis shows that the text introduces a distinct, non-scholastic technical terminology explicitly termed \textit{malla-bhāṣā}. 

The mechanical vocabulary of holds (\textit{lāgas}) did not emerge from clean Sanskrit roots; it was a specialized subaltern argot used within wrestling pit gymnasia (\textit{akhāḍās}) that elite text-composers later forced into formal Sanskrit casing. This linguistic containment is quantitatively verified by our rolling Jaccard indexing. By checking shared localized neighborhoods within a strict proximity boundary ($\pm 5$ words), our matrix highlights a strong context coefficient between the martial subject and vertical column structures (\textit{malla-stambha} = $0.0312$), pinpointing the exact coordinates where tactical somatic skills undergo sovereign regularization.

\subsection{Intertextual Synthesis: Frontier Goddess Assimilation and Alchemy}
This cross-sectarian process of structural capture maps directly onto the early geopolitical distribution of the warrior goddess complex along the Vindhya mountain corridors. As \citet{Yokochi1999} establishes through text-critical tracking, the early historical trajectory of the non-orthodox goddess who abides in the Vindhya mountains (\textit{Vindhyavāsinī}) underwent a deliberate fusion with the pan-Indic \textit{Mahisāsuramardinī} mythos between the fourth and seventh centuries CE. 

Both early Vaiṣṇava and Śaiva state apparatuses vied to absorb these fierce, marginal warrior motifs into their respective mythic cycles—culminating in the original \textit{Skandapurāṇa} subordinating the localized tree-goddess under Pārvatī’s domestic mandala \citep{Yokochi1999}. This macro-historical enclosure mirrors the exact trajectory of scholastic capture, where the untamed violence of the frontier body is progressively harnessed to legitimate royal sovereignty. 

This synthesis is further verified by our link data, which exposes a strong overlap between alchemical calcination and vertical stabilization (\textit{stambha-māraṇa} = $0.0367$), illustrating a shared material vocabulary of containment that preceded its late-medieval philosophical sterilization.

% =====================================================================
% CONSOLIDATED SECTION 6: THE DIACHRONIC CLIMAX AND STATIC ENCLOSURE
% =====================================================================
\section{The Diachronic Climax: Theological Cleaning and Static Enclosure}

\subsection{The Shifting Semantics of Sthāna and Āsana}
The terminal phase of postural flattening operates through a subtle but total shift in the lexicographical boundaries of somatic manuals. Throughout the early medieval strata, terms like stance (\textit{sthāna}) and posture (\textit{āsana}) defined volatile, dynamic positions designed for immediate battlefield utility, close-quarters submission wrestling, or high-line apparatus weapon stabilization. 

As these assets were captured by late-medieval monastic redactors, the physical vocabulary was systematically detached from its material environments. Moving vectors were turned inward, and active defensive holds were re-coded as static, meditative enclosures optimized for internal energy circulation and psychological containment. This linguistic sanitization stripped the kinetic subject of its tactical, open-ended mobility, turning the body into a safe, predictable monument of orthodox theology.

\subsection{Fluid Management: Internalization of the Martial Lock}
This process of structural containment is deeply reflected in how physical locks (\textit{bandhas}) and seals (\textit{mudrās}) were internalized. The mechanical locks originally engineered within the wrestling pit (\textit{akhāḍā}) to break joints, restrict an opponent's blood flow, or absorb high-impact trauma using external wooden levers were systematically decoupled from close-quarters combat. 

Monastic text-composers re-coded these physical techniques as internal fluid management systems, claiming that locking the throat (\textit{jālandharabandha}) or sealing the perineum (\textit{mūlabandha}) was meant to trap the internal nectar (\textit{amṛta}) and prevent its consumption by the gastric fire. This transition demonstrates a clear axis of elite capture: the raw physical survival tactics of marginal groups were stripped of their combat realities and locked inside a passive, metaphysical system of bodily maintenance.

\subsection{Colonial Degradation of the Regional Military Labor Market}
The structural degradation of these indigenous combat gymnasia was deeply accelerated by colonial shifts in state organization. Throughout the early modern period, wrestling systems and vertical apparatus training operated within an active, highly fluid regional ``military labor market,'' offering martial specialists clear paths to political employment, courtly patronage, and localized state defense \citep{Thomas2025}. 

Under British colonial centralization and the enforcement of strict anti-martial regulations, this functional utility was completely severed. Stripped of imperial court funding, the functional mechanics of combat-ready posture in North Gujarat were permanently reduced to a demeaning, token street performance layout: entering the wedding processions of Leuva Patidar elites to execute a single, symbolic wrestling stance (\textit{chapiye dāvu}) in exchange for meager monetary tips (\textit{jamaṇ}) \citep{Sandesara1948}. 

This terminal degradation, explicitly banned by the Jyeṣṭhīmalla caste association in 1945 to safeguard their remaining social dignity, perfectly illustrates the breakdown of indigenous martial capital under the pressures of colonial pacification.

\subsection{Muscular Nationalism and Bourgeois Spectacular Display}
The ultimate enclosure of the rogue body within respectable state structures found its ideal ideological vehicle in the late nineteenth- and early twentieth-century physical culture movements. The printing footprint of the Balodyan Press in 1929—exemplified by Khāsgīvāle’s regional manual \textit{Vet Āṇi Kustī}—marks the precise historical moment where the fluid, dangerous, and criminalized border-evasion skills of the subaltern underbelly were permanently re-coded into a structured, pedagogical syllabus of national pride and preventative healthcare \citep{Khasgivale1929}. 

As further demonstrated by our computational link data, the high contextual overlap between the stabilization pole and zoomorphic ritual conduits (\textit{stambha-vāhana} = $0.0494$) anchors this transition, highlighting how ancient, localized tribal emblems (\textit{vāhanas}) were systematically regularized into a unified national display infrastructure \citep{Abbasi2001}. 

By isolating these complex physical postures from their historical contexts of shadow intelligence gathering and defensive statecraft, these print networks finalized the long-term project of theological sanitization. They flattened the human biological container into a harmless, respectably bourgeois display technology (\textit{prekṣaṇīya kām}) optimized for modern civic consumption.

% =====================================================================
% MASTER BIBLIOGRAPHY (COMPLETELY CLEANED & ALPHABETIZED)
% =====================================================================
\begin{thebibliography}{99}

\bibitem[Abbasi and Tiwari, 2001]{Abbasi2001}
Abbasi, A.~A. and S.~K. Tiwari (2001). ``Contemporary Tribals and Their Totems.'' In \textit{Dimensions of Human Cultures in Central India: Professor S. K. Tiwari Felicitation Volume}, edited by A.~A. Abbasi, 206. New Delhi: Sarup \& Sons.

\bibitem[Birch, 2024]{Birch2024}
Birch, Jason (2024). \textit{The Amaraugha and Amaraughaprabodha of Gorakṣanātha: A Critical Edition with Translation and Study}. Pondicherry: Institut Français de Pondichéry / École française d'Extrême-Orient.

\bibitem[B\"{u}hler, 1875]{Buhler1875}
Bilha\textsubscript{n}a, Vidy\={a}pati (1875). \textit{Vikram\={a}\textsubscript{n}kadevacaritam: A Life of King Vikram\={a}ditya Tribhuvanamalla of Kaly\={a}\textsubscript{n}a}. Edited by Georg B\"{u}hler. Bombay: Government Central Book Depot.

\bibitem[Chitgopekar, 2002]{Chitgopekar2002}
Chitgopekar, Nilima (2002). \textit{Invoking Goddesses: Gender Politics in Indian Religion}. Shakti Books. New Delhi: Har-Anand Publications.

\bibitem[Clark, 2006]{Clark2006}
Clark, Matthew (2006). \textit{The Daśanāmī-Saṃnyāsīs: The Integration of Ascetic Lineages into an Order}. Leiden: Brill.

\bibitem[Clark, 2018]{Clark2018}
Clark, Matthew (2018). Akhāṛās: Warrior Ascetics. In K. A. Jacobsen, H. Basu, A. Malinar, \& V. Narayanan (Eds.), \textit{Brill’s Encyclopedia of Hinduism Online}. Leiden: Brill.

{\bibitem[Danta, 2015]{Danta2015}}
Danta, Chris (2015). Acrobats and Ascetics: Peter Sloterdijk and the Aesthetics of Verticality. \textit{European Journal of English Studies}, 19(1), 66--80.

\bibitem[Das, 1968]{Das1968}
Das, Veena (1968). ``A Sociological Approach to the Caste Puranas: A Case Study.'' \textit{Sociological Bulletin}, 17(2), 141--164.

\bibitem[Debnath, 2024]{Debnath2024}
Debnath, Kunal (2024). \textit{Caste, Marginalisation, and Resistance: The Politics of Identity of the Naths (Yogis) of Bengal and Assam}. Leiden and Boston: Brill (Studies in Critical Social Sciences / New Scholarship in Political Economy, Volume 274).

\bibitem[Dehejia, 1986]{Dehejia1986}
Dehejia, Vidya (1986). \textit{Yogin\={\i}, Cult and Temples: A Tantric Tradition}. New Delhi: National Museum.

\bibitem[Einoo, 1999]{Einoo1999}
Einoo, Shingo (1999). ``The Autumn Goddess Festival: Described in the Pur\={a}\textsubscript{n}as.'' In \textit{Living with \'Sakti: Gender, Sexuality and Religion in South Asia}, edited by Masakazu Tanaka and Musashi Tachikawa. \textit{Senri Ethnological Studies}, 50, 33--70. Osaka: National Museum of Ethnology.

\bibitem[Kautilya, 1969]{ArthaŚ}
Kautilya (1969). \textit{Arthaśāstra of Kautilya}. Edited by R. P. Kangle. Bombay: Motilal Banarsidass. (Digital edition integrated via Hellwig's Digital Corpus of Sanskrit node records).

\bibitem[Keśava, 1982]{KauśP}
Keśava (1982). \textit{Kauśikapaddhati of Keśava}. Edited by V. P. Limaye and R. D. Vadekar. Pune: Tilak Maharashtra Vidyapeeth.

\bibitem[Kh\={a}sg\={\={\i}}v\={a}le, 1929]{Khasgivale1929}
Kh\={a}sg\={\={\i}}v\={a}le, Anant Harihara (1929). \textit{Vet \={A}\textsubscript{n}i Kust\={\i} (Mallakh\textsubscript{a}\textsubscript{m}b)}. Pune: Balodyan Press.

\bibitem[Kullūkabhaṭṭa, 2003]{GRETIL_vcp}
Kullūkabhaṭṭa (2003). \textit{Manvarthamuktāvalī (Commentary on the Manusmṛti)}. GRETIL Digital Repository.

\bibitem[Lorenzen, 1978]{Lorenzen1978}
Lorenzen, David N. (1978). Warrior Ascetics in Indian History. \textit{Journal of the American Oriental Society}, 98(1), 61--75.

\bibitem[Lorenzen and Mu\~{n}oz, 2011]{Lorenzen2011}
Lorenzen, David N. and Adrián Muñoz, eds. (2011). \textit{Yogi Heroes and Poets: Histories and Legends of the Nāths}. Albany: State University of New York Press.

\bibitem[Mallinson, 2007]{Mallinson2007}
Mallinson, James (2007). \textit{The Khecar\={\i}vidy\={a} of \={A}d\={\i}n\={a}tha: A Critical Edition and Annotated Translation of an Early Text of Ha\textsubscript{t}hayoga}. Routledge Studies in Tantric Traditions. London and New York: Routledge.

\bibitem[Mallinson and Singleton, 2017]{Mallinson2017}
Mallinson, James and Mark Singleton, trans. and eds. (2017). \textit{Roots of Yoga}. London: Penguin Classics.

\bibitem[McCartney, 2021]{McCartney2021}
McCartney, Patrick S. D. (2021). \textit{The Not So United States of Yogaland: Transnational Wellness Capitalism and the Geopolitics of Erasure}. London: Anthem Press.

\bibitem[McCartney, 2023]{McCartney2023}
McCartney, Patrick S.D. (2023). ``Poles apart? From Wrestling and Mallkhāmb to Pole Yoga." \textit{Journal of Yoga Studies: Yoga and the Physical Practices of South Asia}, 4, 215–270. \url{https://doi.org}

\bibitem[McCartney, 2025]{McCartney2025}
McCartney, Patrick S.~D. (2025). ``Yoga's `Contortionist Turn,' Pole Acrobatics and Agricultural Festivals: Re-upping Royal Power and Re-ordering the Cosmos.'' \textit{Electronic Journal of Vedic Studies}, 30(3), 1--81. \url{https://doi.org}

\bibitem[McCartney, 2026]{McCartney2026a}
McCartney, Patrick S. D. (2026). \textit{Imagining Sanskritland: `Sanskrit-speaking' Villages, Linguistic Utopias and the Metaphysics of Development}. London and New York: Routledge (Routledge Studies in South Asian History).

\bibitem[McCartney, Forthcoming-a]{McCartneyForthcoming-a}
McCartney, Patrick S. D. (2026). \textit{The Acro-Yoga Complex and Ha\d{t}hayoga's `Contortionist Turn': An Ethno-Historical Study of Ritual and South Asia's `Pole Dancing' Acrobats}. \textit{Abhidha}, 5, Forthcoming.

\bibitem[McCartney, Forthcoming-b]{McCartneyForthcoming-b}
McCartney, Patrick S. D. (2026). \textit{Poison, Vision and Ecstasy: A Transcultural Study of Pharmakon and Ritual Power from the Veda to Greek Mystery Traditions}. \textit{Electronic Journal of Vedic Studies}, Forthcoming.

\bibitem[McCartney, Forthcoming-c]{McCartneyForthcoming-c}
McCartney, Patrick S.~D. (forthcoming). ``The Occupational Overlap between Yoga, Acrobatics and the Occult: Conjuring, Necromancy, Sorcery and the `The Contortionist Turn'.'' \textit{International Journal of Hindu Studies}.

\bibitem[Modha Brahmins, 1964]{MP_7}
Modha Brahmins (1964). \textit{Mallapurāṇa: A Medieval Sanskrit Text on Malla-yuddha}. Edited by B. J. Sandesara and R. N. Mehta. Baroda: Oriental Institute (Gaekwad's Oriental Series).

\bibitem[O'Hanlon, 2007]{OHanlon2007}
O'Hanlon, Rosalind (2007). Military Sports and the History of the Martial Body in India. \textit{Journal of the Economic and Social History of the Orient}, 50(4), 490--523.

\bibitem[Olivelle, 1987]{Olivelle1987}
Olivelle, Patrick (1987). King and Ascetic: State Control of Asceticism in the Artha\v{s}\={a}stra. \textit{Adyar Library Bulletin}, 51, 39--59.

\bibitem[Olivelle, 2011]{Olivelle2011}
Olivelle, Patrick (2011). \textit{Ascetics and Brahmins: Studies in Ideologies and Institutions}. London: Anthem Press.

\bibitem[Patwardhan, 1927]{Patwardhan1927}
Patwardhan, N. G. (1927). \textit{Malla-Vidyā va Pehalavānī Kāmat [The Science of Wrestling and Pugilistic Maneuvers]}. Poona: Chitrashala Press.

\bibitem[Pinch, 1996]{Pinch1996}
Pinch, William R. (1996). \textit{Peasants and Monks in British India}. Berkeley: University of California Press.

\bibitem[Pinch, 2012]{Pinch2012}
Pinch, William R. (2012). \textit{Warrior Ascetics and Indian Empires}. Cambridge: Cambridge University Press.

\bibitem[Puranic Composers, 1980]{ViP_5}
Puranic Composers (1980). \textit{The Viṣṇupurāṇa: With the Commentary of Śrīdhara Svāmī}. Edited by J. L. Shastri. Delhi: Motilal Banarsidass.

\bibitem[Rādhākāntadeva, 2010]{DCS_skd}
Rādhākāntadeva (2010). \textit{Śabdakalpadruma Lexicon data layers}. Digital Corpus of Sanskrit (DCS).

\bibitem[Sandesara, 1948]{Sandesara1948}
Sandesara, Bhog\={\i}l\={a}l J. (1948). \textit{Jye\textsubscript{s}\textsubscript{t}h\={\i}malla J\~n\={a}ti ane `Malla Pur\={a}\textsubscript{n}u'}. Amad\={a}v\={a}d: Gujarat Vidya Sabha.

\bibitem[Shastri and Bhatt, 1991]{Shastri1991}
Shastri, J.L. and G.P. Bhatt, eds. (1991). \textit{The Skandapur\={a}\textsubscript{n}a. Part X: K\={a}\'{s}\={\i}-kha\textsubscript{n}\textsubscript{d}a (P\={u}rv\={a}rdha)}. Translated by G.V. Tagare. Ancient Indian Tradition and Mythology Series, Vol. 49. Delhi: Motilal Banarsidass.

\bibitem[Someśvara III, 1939]{Mān_3}
Someśvara III (1939). \textit{Mānasollāsa of King Bhūlōkamalla Someśvara}, Volume II. Edited by G. K. Shorigondekar. Baroda: Oriental Institute (Gaekwad's Oriental Series).

\bibitem[Srinivas, 2012]{Srinivas2012}
Srinivas, Tulasi (2012). Relics of Faith: Fleshly Desires, Ascetic Disciplines and Devotional Affect in the Transnational Sathya Sai Movement. In B. S. Turner (Ed.), \textit{Routledge Handbook of Body Studies} (pp. 185--205). London and New York: Routledge.

\bibitem[Thomas, 2025]{Thomas2025}
Thomas, Ajay Jacob (2025). ``From \textit{Mallapur\={a}\textsubscript{n}a} to \textit{The Science of Wrestling. Volume I}: Changing Facets of Indian Wrestling Manuals.'' \textit{Indian Historical Review}, 52(1), 116--140.

\bibitem[Veda Vyāsa, 1969]{HV_App}
Veda Vyāsa (1969). \textit{The Harivaṃśa: Being the Khila or Supplement to the Mahābhārata}. Edited by P. L. Vaidya. Poona: Bhandarkar Oriental Research Institute.

\bibitem[Wujastyk, 2025]{Wujastyk2025}
Wujastyk, Dagmar, ed. (2025). \textit{Indian Alchemy: Sources and Contexts}. South Asia Research. New York: Oxford University Press.

\bibitem[Yokochi, 1999]{Yokochi1999}
Yokochi, Yoko (1999). ``The Warrior Goddess in the Dev\={\i}m\={a}h\={a}tmya.'' \textit{Journal of Indian and Buddhist Studies}, 37(2), 43--70.

\end{thebibliography}
